% Chapter 3: Graphs Can Speak
\chapter{Graphs Can Speak}

\step{A Counterintuitive Opening}

Consider two records of a hundred-meter final.

The first is a finish-line photograph: runner A's chest crosses first, about half a body length ahead of B. The judges declare A the champion. Nobody objects.

The second is a radar speed record, measuring both runners every few tenths of a second throughout the race. Lay out the data and something odd appears: from three seconds after the starting gun until just before the finish, B's speed is almost always greater than A's. B started slowly but chased hard in the middle of the race and reached a higher top speed.

So who actually ``ran faster'' in this race?

If faster means finishing a hundred meters first, the answer is A. If it means having the greater speed, then for most of the race the answer is B. Both statements are correct, yet they award different champions. When we casually say ``He runs faster than I do,'' we slide between these two meanings without ever separating them.

When physicists encounter one sentence with two meanings, they do not settle it by arguing; they adopt a more precise language. You already know the language this chapter introduces: draw a graph. You will see that graphing the race makes the difference between A and B obvious at a glance, with no argument required. You will also see that a line's steepness and an area on a graph correspond to two of the physical world's most important questions: ``How fast is it changing now?'' and ``How much has accumulated altogether?'' These questions run through the entire book, from motion to forces, from electric current to heat, and all the way to quantum mechanics.

Incidentally, the invention of this language was itself a revolution. Before it, people could describe change only in words---``faster and faster,'' ``gradually slowing down''---without specifying exactly how fast faster was. After it, change could for the first time be discussed precisely, calculated, and drawn.

\step{Make a Guess First}

Before reading on, write down your guesses. At the end of the chapter, return and see how much your intuition got right.

Question one. A transparent cup collects water from a faucet at a constant flow rate. How does the water level change with time? Would a graph of water height against time be a straight line, an upward-bending curve, or a downward-bending curve?

Question two. An elevator leaves the first floor, stops once at the fourth, and finishes at the eighth. On an elevator position--time graph, what does the line look like during the dozen or so seconds it stops? Does it keep rising, or turn a corner?

Question three. A car's dashboard has two gauges: the speedometer shows the current speed, and the odometer the total distance traveled. Suppose you press the accelerator and the speedometer climbs from forty to eighty kilometers per hour. How does the increase in the odometer during that interval relate to how quickly the speedometer needle sweeps upward? Can a speed record give distance traveled, and conversely, can a distance record give speed?

Question four, the most subtle. Someone walks north, turns around, and walks south to the starting point. On a velocity--time graph, should the southbound portion be above or below the horizontal axis? Your intuition may protest that velocity cannot be negative. Remember that reaction; later this chapter will formally overturn it.

There is no need to rush an answer. Next we will let your hands and eyes find out before your mind does.

\step{Hands-on Experiment}

\begin{experiment}[A Water Clock: Draw a Straight Line Yourself]
You need a transparent glass, or a cut-open water bottle; a marker; a ruler; a phone stopwatch; and a dripping faucet or a plastic cup with a small hole.

Step one: adjust the faucet to a small, steady stream that flows into the glass at a constant rate.

Step two: once the flow is steady, start the stopwatch. Every twenty seconds, mark the water level on the side of the glass, making six to eight marks in all. Work quickly and keep timing; do not stop the watch between marks.

Step three: turn off the water. Measure each mark's height above the bottom and record its corresponding stopwatch reading. You now have a two-column table: time and the corresponding water height.

Step four: take graph paper, or draw your own grid. Put time on the horizontal axis and water height on the vertical axis. Plot each pair of numbers as a point, then examine their arrangement.

If the flow was steady enough, the points almost form a straight line. This is one of the chapter's most important images: \textbf{uniform change appears as a straight line}. The steeper the line, the faster the water rises and thus the greater the flow.

Now make a comparison: squeeze the hose to vary the flow from large to small and repeat. This time the points form a broken line that is steep first and flatter later. Without a teacher, you have drawn this chapter's protagonist: a graph that speaks.
\end{experiment}

There is a hidden variation, reserved for the challenges at the chapter's end: if you replace the straight-sided glass with a cup wider at the top than at the bottom, does the water level still rise along a straight line? Do not peek at the answer yet.

\step{The Old Model Runs into Trouble}

You may ask why words, tables, and the average velocity learned in middle school are not enough. Each has a defect it cannot cure.

Start with words. ``The train gets faster after leaving the station, then gradually slows as it approaches the next.'' Everyone understands, but this answers no specific question: how much faster, over how many seconds, and where did the most abrupt change occur? Describing change verbally is like describing a person's height only as ``quite tall'' or ``rather tall.'' Fine for conversation, useless for tailoring a suit.

Next, tables. They are precise, but so honest that they tell you everything and therefore almost nothing. A high-speed train's record with one speed per second contains six hundred numbers in ten minutes. Staring at those numbers, it is hard to see whether the train began slowing in the third minute or the seventh. The pattern hides in numbers, and human eyes are poor at seeing shapes in numerical queues.

Finally, average velocity, the most concealed difficulty. A bus covers ten kilometers in thirty minutes, for an average of twenty kilometers per hour. True, but that sentence flattens every story along the way: zero speed at red lights, low speed after leaving a stop, high speed across the Yangtze River bridge. An average is like describing a river as having an average depth of half a meter. That sentence could kill a swimmer if the middle is three meters deep. Averages answer ``How is it overall?'' while deliberately declining to answer ``How is it now?'' Physics often needs precisely now: speed at the instant of collision determines whether an airbag should deploy, and speed at impact determines whether an egg breaks.

The phrase ``velocity at this instant'' is more troublesome still. Velocity is distance traveled during an interval divided by that interval. What does velocity mean for a single instant, an interval of zero duration? Zero divided by zero? Our old concept hits a wall.

The navigation app you use daily relies heavily on the average-speed model. Its forecast of forty minutes to arrival divides the remaining distance by estimated average speeds on the remaining road sections. When traffic ahead jams and alternates between fast and slow, the forecast falters and the arrival time jumps later and later. The app has not become stupid: a model that hides the present inside an average naturally fails when conditions vary sharply. It gives a statement about averages while you want the actual situation at each moment.

Historically, people were stuck at this wall for nearly two thousand years. The ancient Greek Zeno mocked it with his arrow paradox: an arrow occupies a definite position at every instant, so how can it move? Tearing down this wall requires a new conceptual apparatus, not cleverer sophistry.

\step{Inventing New Concepts}

The apparatus has three layers. Let us build them one at a time.

Layer one: \textbf{coordinates}. To discuss where something is, first name every position. Lay a graduated axis along a road, choose a point as zero, and designate a positive direction. Each position then receives a number: three meters right of zero is \(+3\), two meters left is \(-2\). The minus sign does not mean less than nothing; it means the other side of zero. This foreshadows the negative velocity in question four: signs indicate direction, not quantity. This apparently ordinary step is bold: it declares that positions can be recorded numerically.

One easily overlooked point deserves emphasis: you are entirely free to choose the origin and positive direction. Set the track's finish line as zero and west as positive, and all the numbers change. Yet the graph's shape, the sign relationships between slopes, and differences between areas---all the physics---remain unchanged. Coordinates are labels we stick on the world, not part of the world itself. The modest-looking idea that physical laws do not depend on how we label things will grow into one of the book's central pillars when we discuss reference frames and relativity.

Layer two: \textbf{variables and functions}. Every row in the water-clock record is a pair of numbers: a time and a height. If the rule is definite---give me a time and I give you the height then---we say height is a function of time. What kind of machine is a function? Like a vending machine, it takes one number and returns one definite number, no more or less, without regard to its mood. That is the resemblance. The difference matters too: the machine does not literally contain heights, and a function is not a scroll concealed inside an object. It is only our ledger for recording and predicting a correspondence. Moreover, a functional relationship need not be causal: ``At three in the afternoon the temperature is twenty degrees'' does not mean three o'clock caused twenty degrees. We will return to this when discussing experimental data.

Layer three: \textbf{graphs}. Each number pair in the ledger becomes a point on paper: measure time horizontally and height vertically. Connect the points and you have a graph. Its magic is compressing the entire ledger into a shape you can see at a glance. Six hundred train-speed numbers reveal little, but six hundred connected points immediately show the minute when slowing began. What words cannot explain and tables bury, a graph announces through shape.

The shapes contain two codes, the real protagonists of this chapter.

Code one: \textbf{slope answers ``How fast is it changing?''} Take two points on the graph and divide their vertical difference by their horizontal difference: that is slope. On a position--time graph, the vertical difference is change in position and the horizontal difference is change in time. Slope therefore means position change per unit time, precisely velocity. Steeper means faster change; a horizontal line has zero slope and means no change; a descending line has negative slope and means change in the negative direction. Recall the water clock: constant flow gives a uniformly rising water level, a straight graph, and the same slope everywhere. Flow that starts large and becomes small produces a graph that starts steep and flattens, with a slope that decreases. Slope is the precise version of steepness.

What is slope like? Like a hillside's steepness: height climbed divided by horizontal distance advanced. That is the resemblance. There are two differences. First, real hillsides are constrained by the soil, while a graph's slope can become arbitrarily large. As a line becomes vertical, its slope tends to infinity, physically meaning change with no elapsed time, usually a warning that the model is wrong. Second, a hillside's steepness is what your eye sees, but a graph's apparent steepness depends on axis scales. Stretch the horizontal axis and the same data look gentler. Slopes must therefore be reported with units, such as three meters per second. Saying only ``very steep'' means nothing.

Code two: \textbf{area answers ``How much has accumulated?''} Change perspective. If you know the velocity at every moment and want to know how far the object went, what do you do? Divide time into many short intervals, each short enough that velocity barely changes. Displacement in each is roughly velocity times interval duration. On a velocity--time graph, that product is the area of a thin rectangle, with velocity as height and time as width. Adding all the small rectangles gives the total area between the graph and the horizontal axis. Thus, \textbf{the area under a velocity--time graph is the total change in position during the interval}. At constant velocity the line is horizontal, the area a rectangle: the elementary formula ``speed times time equals distance'' is a special case of area.

What is area like? Like a water-meter reading: the meter does not directly tell you the present flow, but quietly adds the flow at every moment. That is the resemblance. The difference is that a graph's area is not square centimeters measured with a ruler. Its units are vertical-axis units multiplied by horizontal-axis units: velocity times time gives meters, not square meters. Geometry is merely the shape it borrows.

An even more elegant correspondence deserves remembering separately. Slope and area are inverse operations: measure slopes on a position graph to get velocity, and measure areas on a velocity graph to recover position. A car's dashboard puts this correspondence before your eyes: the speedometer is a slope gauge, the odometer an area gauge, and the two always agree. This is the central idea of calculus, which the mathematical bridge will discuss.

\step{The Model in One Sentence}

\begin{oneline}
On a position--time graph, slope tells you how fast things are changing now; under a velocity--time graph, area tells you how much has accumulated altogether---slope handles rates, area handles totals, and the operations are inverses.
\end{oneline}

Scan the world with this sentence and you will find it everywhere. Water and electricity meters record area, accumulation; water flow and the operating state of an electric water heater are slopes, rates. In phone health apps, current pace is slope and today's step count converted into distance is area. In a bank account, the monthly rate of salary flowing in and expenses flowing out is slope, while the balance is area. Many branches of physics amount to using these two tools repeatedly in different settings.

A further reminder: this language is not exclusive to motion. A hospital electrocardiogram has voltage vertically and time horizontally; a doctor reads the shape of voltage changes. A weather forecast's temperature curve has positive slope on a warming morning and negative slope on a cooling evening. The area under a rice cooker's power--time graph is the electrical energy consumed in that cooking session, the source of the kilowatt-hours on your bill. Graphs apply whenever one quantity changes with another. This chapter uses motion simply because it is easiest to see and touch.

\step{The Mathematical Translator}

Now translate the intuitions into symbols one at a time. Remember the order: first a picture, then a formula.

First translation: slope. On a position--time graph, choose times \(t_1\) and \(t_2\), with corresponding positions \(x_1\) and \(x_2\). The average velocity over that interval is the slope of the line joining the two points:
\[
\bar{v}=\frac{\Delta x}{\Delta t}=\frac{x_2-x_1}{t_2-t_1}.
\]
Check it as usual. If position is unchanged, \(\Delta x=0\), the average velocity is zero: a stationary person indeed has zero average velocity. For the same displacement in less time, the smaller denominator gives a larger value: covering the same distance more quickly does mean moving faster. Turn around and walk back, with \(x_2<x_1\), and \(\Delta x\) and average velocity are negative: the sign faithfully reports reversal. This answers the fourth opening question: the southbound portion of the velocity--time graph belongs below the horizontal axis.

\begin{center}
\begin{tikzpicture}[>=Stealth,scale=1.05]
  \draw[->] (0,0) -- (5.6,0) node[below left] {\(t\) (time)};
  \draw[->] (0,0) -- (0,3.4) node[below left] {\(x\) (position)};
  \draw[thick] (0,0.4) -- (4.8,3.0);
  \draw[dashed] (1.5,1.19) -- (3.8,1.19) node[midway,below] {\(\Delta t\)};
  \draw[dashed] (3.8,1.19) -- (3.8,2.50) node[midway,right] {\(\Delta x\)};
  \fill (1.5,1.19) circle (1.2pt) (3.8,2.50) circle (1.2pt);
  \node[align=center] at (1.9,2.9) {Slope \(\displaystyle=\frac{\Delta x}{\Delta t}\)};
\end{tikzpicture}
\end{center}

Second translation: from average to instantaneous. Average velocity flattens an interval of time. The remedy is direct: shorten the interval. An average over one second comes closer to now than one over ten seconds, and an average over a tenth of a second closer still. Continue shortening, and the values approach a definite number steadily. We define that number as the \textbf{instantaneous velocity} at that time. Intuitively, it is the slope of the graph's tangent at that point. Imagine magnifying a tiny stretch repeatedly: the curve looks more and more like a straight line, whose slope is the instantaneous velocity.

Notice what this definition does. It never calculates zero divided by zero. It calculates a sequence of average velocities over progressively shorter intervals and asks where the sequence tends. That is how Zeno's arrow is acquitted: velocity at an instant is no longer a contradiction, but shorthand for a limiting process. There is no mysticism, only a convention: as intervals shrink and readings settle to an unchanging value, we call that value instantaneous velocity. Real instruments can only use finite intervals, of course, but this does not invalidate the clean definition, any more than your inability to draw a perfect circle spoils the concept of a circle.

Third translation: area. Given velocity \(v\) at every instant, find the position change from \(t_1\) to \(t_2\). Divide time into short intervals of width \(\Delta t\), with nearly constant velocity in each. Each displacement is approximately \(v\,\Delta t\); add them:
\[
\Delta x\approx v_1\,\Delta t+v_2\,\Delta t+v_3\,\Delta t+\cdots
\]
Finer intervals give a better approximation. In the limit, the sum is the area under the velocity--time graph.

\begin{center}
\begin{tikzpicture}[>=Stealth,scale=1.05]
  \draw[->] (0,0) -- (5.6,0) node[below left] {\(t\) (time)};
  \draw[->] (0,0) -- (0,3.4) node[below left] {\(v\) (velocity)};
  \fill[blue!15] (0,0) -- (0,0.8) -- (4.2,2.9) -- (4.2,0) -- cycle;
  \draw[thick] (0,0.8) -- (4.2,2.9);
  \draw[dashed] (4.2,2.9) -- (4.2,0);
  \node[below] at (4.2,0) {\(t_2\)};
  \node at (1.9,1.05) {Area \(=\) change in position};
\end{tikzpicture}
\end{center}

Check as usual. Zero velocity puts the graph on the horizontal axis, with zero area and no change in position: the odometer does not move when parked. Constant velocity \(v\) gives a horizontal line and a rectangle of area \(v(t_2-t_1)\), recovering the elementary speed-times-time formula. Negative velocity places the line below the axis; by convention its area is negative and position change is negative too. When the car reverses, its displacement ledger decreases. We mean a displacement ledger here, not a broken odometer; this detail will return at the chapter's end.

The fourth translation combines the first two and supplies a real engineering tool. A railway train diagram is essentially a position--time graph: time runs horizontally and stations along the line vertically. Each train is a sloping line, with slope equal to speed. Two intersecting lines mean two trains at the same place and time---a collision. In graphical language, a dispatcher's entire task is to keep hundreds of lines from intersecting. A fast train catching a slower one in the same direction appears as a steep line approaching a gentler one. The dispatcher must divert the slow train to a siding to stop briefly: on the graph, add a horizontal segment to the gentle line until the steeper one passes. Behind the announcement ``Stopping temporarily to allow another train to pass'' lies this graph.

Now use real numbers to see how handy area can be. A high-speed train takes roughly five minutes to accelerate from rest to 350 kilometers per hour, about \(97\) meters per second. Its velocity--time graph during acceleration is approximately a triangle, with a time base of \(300\) seconds and a height of \(97\) meters per second. The distance traveled while accelerating is its area:
\[
\Delta x\approx\frac{1}{2}\times 97\ \text{m/s}\times 300\ \text{s}\approx 1.5\times10^{4}\ \text{m},
\]
or about fifteen kilometers. Without solving any equation, just measuring an area, we know that a train has quietly traveled more than ten kilometers between starting and reaching full speed. That is why high-speed railways need such long, level, straight stretches.

\begin{mathbridge}
If you know a little function notation, the limit above can be written more compactly. Instantaneous velocity is
\[
v(t)=\lim_{\Delta t\to 0}\frac{x(t+\Delta t)-x(t)}{\Delta t}=\frac{dx}{dt},
\]
read as the derivative of \(x\) with respect to \(t\). It is a function, assigning a slope to each instant. For a concrete example, let \(x(t)=2t^{2}\), with position in meters and time in seconds. Then
\[
\frac{x(t+\Delta t)-x(t)}{\Delta t}=\frac{2(t+\Delta t)^{2}-2t^{2}}{\Delta t}=4t+2\,\Delta t.
\]
As \(\Delta t\) tends to zero, the \(2\,\Delta t\) term disappears, leaving \(v(t)=4t\). Check dimensions using Chapter 2's tool: in \(2t^{2}\), the coefficient \(2\) actually carries units of meters per second squared, so \(4t\) has units of meters per second and is indeed a velocity.

Conversely, area in the limit is written as an integral:
\[
\Delta x=\int_{t_1}^{t_2} v(t)\,dt.
\]
The inverse relationship---differentiate position to obtain velocity and integrate velocity to recover position---has a grand name: the fundamental theorem of calculus. In this chapter's language, \textbf{the area under the slope graph equals the height difference on the original graph}. That graphical sentence is the foundation of the entire edifice of calculus.
\end{mathbridge}

\begin{challenge}
Using area on a velocity--time graph gives a famous formula for free. For uniformly accelerated motion, the velocity graph is a straight line through the origin, \(v=at\). Its area up to time \(t\) is a triangle:
\[
\Delta x=\frac{1}{2}\times t\times at=\frac{1}{2}at^{2}.
\]
The uniform-acceleration displacement formula learned in middle school is simply a triangle's area. There is no need to memorize it.

Go a level deeper: the slope operation can be repeated. Taking slopes on a position graph gives velocity; taking slopes again gives a graph of how quickly velocity changes. This quantity, acceleration, makes its formal entrance in the next chapter. Conversely, areas under an acceleration graph give velocity, and another area calculation gives position. Your phone works this way: its accelerometer measures acceleration hundreds of times a second, and its chip accumulates areas twice to estimate velocity and displacement. Engineering adds a cruel detail: every measurement contains a small error, and area accumulation makes errors grow. Integrate once and the error grows linearly; integrate twice and it grows quadratically. Position estimated solely from an accelerometer therefore drifts wildly after a few minutes, requiring continual resetting with satellite positioning and wireless signals. Slope and area are exact inverses in mathematics, but behave asymmetrically in the noisy real world: taking slopes amplifies noise, while taking areas accumulates bias. This is an essential lesson when graphical language enters engineering.
\end{challenge}

\step{The Model's Limits}

Graphical language is extremely powerful, and therefore extremely easy to misuse. Each of these four boundaries corresponds to a common class of mistakes.

Boundary one: \textbf{a graph is a ledger, not a map}. A position--time graph describes where something is at each time, not the route it follows. An object moving back and forth on a straight line can produce a rising and falling curve. An object actually following a curved path can produce a straight position--time graph if its position along that path changes uniformly. Mistaking the graph's shape for the motion's trajectory is the beginner's most common misreading. A curved graph means the rate of change is changing; it says nothing about turning a corner. Direction needs the vectors introduced in Chapter 4.

Boundary two: \textbf{looking steep does not mean having a large slope}. Stretch the horizontal axis for the same data and the graph looks gentler; compress the vertical axis and it looks steeper. Many newspaper and advertising graphs of soaring performance or plunging results manufacture their visual effect by truncating vertical axes and stretching horizontal ones. On a physics graph, trust only a numerical slope with units, not your visual impression. For every graph, first check which quantities and units the two axes represent, and whether they start at zero or are cut off partway through.

Boundary three: \textbf{area is a signed ledger}. Parts of a velocity--time graph below the horizontal axis count negatively. The area therefore calculates net change in position, displacement, not total distance traveled. Walk ten meters forward and ten back: displacement is zero, distance twenty meters, because the positive and negative velocity areas cancel. To find distance, take the absolute value of velocity before calculating area. This distinction will take center stage when describing motion in the next chapter.

Boundary four: \textbf{a line between two points makes a claim}. Experiments always give discrete points. Joining them into a smooth curve claims that behavior at intermediate times also varies smoothly between them. This is usually right, but is not given for free. During the minute when water boils, its temperature--time graph is flat; melting ice likewise gives a temperature plateau. If measurement points are too sparse, you can miss a plateau entirely and draw a phase transition as a gentle slope. How densely to sample depends on how quickly the process of interest changes. That is part of experimental design, not something plotting software should decide for you.

\step{Explaining the Opening Phenomenon}

Return to the hundred-meter final and plot A and B's positions against time on the same graph.

\begin{center}
\begin{tikzpicture}[>=Stealth,scale=1.05]
  \draw[->] (0,0) -- (5.8,0) node[below left] {\(t\) (time)};
  \draw[->] (0,0) -- (0,3.6) node[below left] {\(x\) (position)};
  \draw[dashed] (0,3) node[left] {\(100\) m} -- (5.2,3);
  \draw[thick] (0,0) -- (3.4,3) node[pos=0.75,above left] {A};
  \draw[thick] (0,0) .. controls (1.6,0.5) and (2.8,2.5) .. (4.4,3) node[pos=0.8,below right] {B};
  \draw[dashed] (3.4,3) -- (3.4,0) node[below] {\(t_{\text{A}}\)};
  \draw[dashed] (4.4,3) -- (4.4,0) node[below] {\(t_{\text{B}}\)};
\end{tikzpicture}
\end{center}

A's graph is almost a steep straight line: a quick start and steady rhythm throughout. B's starts gently, with slow starting reaction and small slope, then grows steeper in the middle until its slope exceeds A's. That is the radar's statement that B was faster mid-race. The two lines reach the hundred-meter height in sequence, A first, B later. Every part of the race---who started quickly, who finished strongly, who crossed first---is visible at a glance.

The opening contradiction dissolves. A velocity--time graph measures how fast position changes now, B's advantage. Race rules care about who accumulates a hundred meters first, whose position changes by a hundred meters sooner: the language of area. B's problem is that the period of greater speed is too short; the race ends before the lost area is recovered. In the one-sentence model's terms, \textbf{B wins on slope, A wins on area}. ``Who is faster?'' sounds ambiguous because everyday language lumps slope and area together under one word. When graphs speak, they separate those two meanings onto two axes.

That fulfills the chapter's opening promise: when one sentence has two meanings, draw a graph instead of arguing.

Test the tools with real data. At the 2009 World Championships in Berlin, Usain Bolt ran a hundred meters in 9.58 seconds. The race's area ledger records a hundred-meter change in position. His average speed was about \(100\div 9.58\approx 10.4\) meters per second, while split times show a top speed around twelve meters per second between sixty and eighty meters. His velocity--time graph therefore climbs rapidly from zero, fluctuates slowly at a high level, and has a total area of exactly a hundred meters. A runner maintaining a constant ten meters per second would have a horizontal velocity graph and rectangular area, completing the race in exactly ten seconds, slower than Bolt. Those 0.42 seconds are hidden in the different graph shapes: Bolt accumulated enough area sooner. A hundred-meter race rewards whoever accumulates a hundred meters of displacement first, not whoever has the biggest speed number. That sentence now has a precise graphical meaning.

\step{Thought Experiments and Challenges}

\begin{thought}[Your Life Is a Line]
Push these tools to an extreme scale. Imagine a huge graph with time horizontally and your position within a city vertically. Draw backward from today: every departure, trip to school, and return home is a segment. A strange perspective appears: you never disappear from this graph. Your life is a continuous line running from left to right. Even if you lie in bed all day, the line does not stop; it simply becomes horizontal, extending quietly deeper into time. Physicists have a name for it: a worldline.

Now consider two extremes. First, what worldline is impossible? One that turns back to the left, putting two versions of you at the same instant and implying travel into the past. Second, what does a perfectly vertical line mean? Position changes while time does not pass: movement over a distance in no time, with infinite slope. Every real object's worldline is constrained to a certain steepness. This is not an engineering limitation but a basic rule of the universe: there is a maximum slope no worldline can exceed, corresponding to the speed of light. When you reach relativity, you will discover that its most important diagram, the spacetime diagram, is essentially this chapter's position--time graph with the axes' roles rearranged. Grand terms such as causality, simultaneity, and time dilation become geometric questions about slopes and angles. The graph-reading skills you have learned here will apply directly.
\end{thought}

\begin{puzzles}
\item An elevator rises from the first floor, stops at the fourth with its doors open for ten seconds, continues to the eighth, and then returns empty to the first. Sketch its position--time graph qualitatively. How does the downward portion's slope relate to the upward portion's slope? \textbf{Hint}: what shape appears while position is constant during the stop? When descending reverses the direction of position change, what happens to the slope's sign? What determines whether its magnitude is greater or smaller going down than going up?

\item In the water-clock experiment, replace the straight-sided glass with a cup wider at the top than at the bottom, such as an ordinary tapered drinking cup, while keeping water flow constant. Is the water-height--time graph still straight? If not, which way does it bend? \textbf{Hint}: slope means the rate at which the water level rises, whereas constant flow means constant volume added per unit time. As the cup gets wider upward, does the same volume raise the surface more or less? Distinguish accumulated volume from changing height: once again, area and slope are different things.

\item Cars A and B travel along the same straight road. Their velocity--time graphs are both straight: A's extends horizontally from sixty kilometers per hour, while B's rises steadily from zero and eventually intersects it. Does the intersection mean B has caught A? \textbf{Hint}: at the intersection, the lines have the same height. What quantity does height represent on this graph? Catching up concerns a different quantity, the difference between the areas under their respective graphs. Whose area is larger at that instant?

\item An elastic ball falls, bounces, falls again, and keeps bouncing to successively lower heights. Sketch its velocity--time graph qualitatively, taking upward as positive, and pay special attention to the instant of ground contact. \textbf{Hint}: falling velocity points downward, is negative, and becomes more negative. After a bounce it points upward and is positive. Contact is extremely brief, so velocity jumps sharply from negative to positive and the graph is nearly vertical, with an enormous slope at that instant. In reality a steeper jump means shorter contact time and a more abrupt velocity change. How abrupt that change is will be the business of force in later chapters.
\end{puzzles}

None of these problems requires a numerical answer, but each tests the same skill: can you let a graph think for you? Slope answers how fast things change, area answers how much accumulates, and signs report direction. With these three tools, the next chapter will give change itself a direction as we enter the world of vectors.
